\section{Empirical reporting contract}

The matched evaluation must be legible as a scientific-process and AI-systems result. The primary question is not whether \RAKL{} uses fewer tokens in isolation, but whether it changes the frontier between valid scientific success and the resources consumed to obtain it. A run counts as a \emph{valid scientific success} only when the final task outcome is correct and all preregistered hard scientific-process gates pass. These include, where applicable, absence of unsupported authority upgrades, mandatory evidence retention, context/estimand alignment, negative-history preservation, correct handling of hidden scientific defects and valid evaluator/corpus identity.

Provider-reported input and output tokens are summed over all model calls in an arm, including planning, retrieval, criticism, recovery and verification. Cached input is separated when exposed by the provider; hidden chain-of-thought counts that are not exposed are not invented. Retrieval/tool calls, retrieved volume, preprocessing, provider cost and wall time remain separate coordinates. For a matched stratum $s$, a headline resource quantity is
\[
\operatorname{TPVS}(s)=
\frac{\sum_{i\in s}\text{billable tokens}_i}
{\sum_{i\in s}\mathbf 1[\text{valid scientific success}_i]}.
\]
Zero valid successes therefore yield an infinite rather than missing cost-per-success value. When several budget levels are run, token reduction is additionally reported at a capability target frozen before outcomes are inspected.

The main empirical figures are organized by reader question. A reliability--cost frontier plots valid scientific success against billable tokens, provider cost and wall time. A process-risk figure reports task-level effect sizes for hidden-defect misses, unsupported authority upgrades, mandatory-evidence omission, failed counterevidence uptake, negative-history loss and false/missed blocks. An architecture-by-evidence-access interaction figure separates gains from better control architecture from gains due merely to public, curated or complete-sealed evidence access. A prospective OWMD figure reports valid hidden-mechanism recall and false discovery versus a frozen search/tool budget, including hidden-name and lexical-distance strata.

All empirical aggregates retain numerator/denominator counts and task-level uncertainty. Repeated generations within a task are nested replicates rather than independent samples. Confirmatory comparisons are paired by task, base model/version, evidence-access arm, resource ceiling, evaluator and seed schedule, with effect sizes and 95\% intervals reported alongside any registered inferential test. The manuscript remains vector-valued: a token saving accompanied by more authority leakage or scientific-defect misses is not a positive result.

Quantitative figures use editable vector text, final-size uncertainty definitions and exact source receipts. They contain no arrows pointing to data, no opaque text boxes over the plotting region and no point-by-point prose annotations. Panel letters remain outside the data region; method names are carried by axes or a shared legend. Per-model replications, token decomposition, latency tails, calibration, ablations and threshold sensitivity are reserved for Extended Data unless an outcome makes them load bearing.

The full metric, figure and QA contract is frozen before outcome inspection in \path{research/PAPER2_EMPIRICAL_FIGURE_CONTRACT_20260810.md}. Result receipts conform to \path{schemas/paper2-empirical-result.schema.json}.